\section{Slope stability and symmetries of the rays}\label{sec:stability}

Let $X=X_\Sigma$ be a smooth toric Fano variety of dimension $n$, with
lattice $N$, dual lattice $M$, and primitive ray generators
$u_\rho\in N$ for $\rho\in\mathcal R:=\Sigma(1)$. Write $D_\rho$ for the
corresponding invariant prime divisor, and put
\[
 d_\rho=D_\rho\cdot(-K_X)^{n-1},\qquad
 D=(-K_X)^n=\sum_{\rho\in\mathcal R}d_\rho,\qquad
 \delta_\rho=\frac{n d_\rho}{D}.
\]
Thus $d_\rho>0$, $\delta_\rho>0$, and $\sum_\rho\delta_\rho=n$.
For a torsion-free sheaf $\mathcal E$, its slope is
$\mu(\mathcal E)=c_1(\mathcal E)\cdot(-K_X)^{n-1}/\operatorname{rk}\mathcal E$.
Stability, respectively semistability, means that every nonzero subsheaf
of smaller rank has strictly smaller slope, respectively no larger
slope. In particular, $\mu(T_X)=D/n$.

\begin{lemma}[Toric slope criterion]\label{lem:klyachko}
The tangent bundle $T_X$ is stable, respectively semistable, if and only
if every nonzero proper complex subspace $V\subset N_\CC$ satisfies
\begin{equation}\label{eq:klyachko}
 w(V):=\sum_{\rho:\,u_\rho\in V}\delta_\rho
 <\dim V,\qquad\text{respectively }w(V)\le\dim V.
\end{equation}
It suffices to test subspaces spanned by ray generators.
\end{lemma}

\begin{proof}
This is the criterion of Hering--Nill--S\"u\ss{}
\cite[Proposition 1.2 and Remark 2.6]{HeringNillSuess2022}, obtained from
Klyachko's filtrations \cite{Klyachko1990}. For completeness, the
decreasing filtrations for $T_X$ on $N_\CC$ are
\[
 E^\rho(i)=
 \begin{cases}
 N_\CC,&i\le0,\\
 \CC u_\rho,&i=1,\\
 0,&i\ge2.
 \end{cases}
\]
The saturated equivariant subsheaf $\mathcal E_V$ associated to $V$ has
induced filtrations $V\cap E^\rho(i)$
\cite[Corollary 0.0.2]{DasguptaDeyKhan2021}, rank $\dim V$, and degree
$\sum_{u_\rho\in V}d_\rho$. Both slope stability and semistability can be
tested on these subsheaves
\cite[Proposition 2.3 and Remark 2.4]{HeringNillSuess2022}.
If $V$ contains rays, their span $V'$ contains exactly the same rays and
has dimension at most $\dim V$, so it has at least the same slope.
If $V$ contains no ray, then $w(V)=0<\dim V$.
\end{proof}

For a rational ray-generated subspace $V$, the saturated subsheaf
$\mathcal E_V$ is also the toric foliation whose general leaves are
orbits of the corresponding subtorus. Its first Chern class is
\[
 c_1(\mathcal E_V)=\sum_{u_\rho\in V}D_\rho.
\]
This follows from the same induced filtrations. In particular, a line
containing a single ray has first Chern class $D_\rho$; a line containing
two opposite rays has the sum of the two divisors.

There is an equivalent convex-geometric normalization. Let
\[
 P=\{m\in M_\RR:\langle m,u_\rho\rangle\ge-1
                    \text{ for all }\rho\}
\]
be the anticanonical polytope, with facets $F_\rho$. With lattice volume,
$H^n=n!\operatorname{Vol}(P)$ and
$D_\rho H^{n-1}=(n-1)!\operatorname{Vol}(F_\rho)$, where $H=-K_X$.
The probability measure assigning to $u_\rho$ the mass
\[
 \nu_P(\rho)=\frac{\operatorname{Vol}(\operatorname{conv}(0,F_\rho))}
                  {\operatorname{Vol}(P)}
            =\frac{\delta_\rho}{n}
\]
is the cone-volume measure, written on ray directions. The criterion
is precisely $\nu_P(V)<\dim(V)/n$. This connection with subspace
concentration is established in \cite{HeringNillSuess2022}; our use of
it will be through the explicit divisor degrees of the two constructions.

\begin{lemma}\label{lem:degreerelation}
The divisor degrees satisfy
\[
 \sum_{\rho\in\mathcal R}d_\rho u_\rho=0.
\]
\end{lemma}

\begin{proof}
For $m\in M$, the principal divisor of the character $\chi^m$ is
$\sum_\rho\langle m,u_\rho\rangle D_\rho$. Its intersection with
$(-K_X)^{n-1}$ vanishes. Pairing with all $m\in M$ gives the assertion.
\end{proof}

\subsection{The ray matroid and invariant subspaces}

The \emph{ray matroid} is the linear matroid on $\mathcal R$ whose rank
function is
\[
 r(S)=\dim_\CC\operatorname{span}\{u_\rho:\rho\in S\},
 \qquad S\subseteq\mathcal R.
\]
A circuit is a minimally linearly dependent set of ray generators.
The components of the matroid are the equivalence classes generated by
requiring all members of a circuit to belong to the same class; a ray
in no circuit forms a singleton class. The matroid is \emph{connected}
if it has a single component.

\begin{lemma}\label{lem:raycomponents}
If $C_1,\ldots,C_s$ are the components of the ray matroid and
$V_a=\operatorname{span}\{u_\rho:\rho\in C_a\}$, then
\[
 N_\CC=V_1\oplus\cdots\oplus V_s.
\]
Every direct-sum decomposition of $N_\CC$ in which each ray generator
lies in a summand is obtained by grouping these component subspaces.
In particular, stability of $T_X$ implies connectedness of its ray matroid.
\end{lemma}

\begin{proof}
Choose a basis $B\subset\mathcal R$. For $\rho\notin B$, the ray $u_\rho$
and the basis vectors occurring with nonzero coefficient in its
expansion form a circuit. All those basis vectors therefore belong to
the same component as $\rho$. Consequently $V_a$ is spanned by the
subset $B\cap C_a$ of this basis, proving the direct sum.
Conversely, a circuit cannot meet two summands of any decomposition
containing every ray in a summand: projecting its dependence to either
summand would give a dependence on a proper subset of the circuit.
Thus each component is contained in one summand. Since the rays span
$N_\CC$, each summand is the sum of the component subspaces it contains.

If $s>1$, the equalities $\sum_a w(V_a)=n=\sum_a\dim V_a$ imply
$w(V_a)\ge\dim V_a$ for at least one component, contradicting
Lemma~\ref{lem:klyachko}.
\end{proof}

\begin{proposition}[Symmetry reduction]\label{prop:symmetry}
Let $A$ be any subgroup of the lattice automorphisms preserving
$\Sigma$. Then $T_X$ is stable if and only if its ray matroid is
connected and
\[
 w(V)<\dim V
\]
for every nonzero proper $A$-invariant subspace $V\subset N_\CC$
spanned by ray generators.
\end{proposition}

\begin{proof}
The group $A$ permutes the rays and preserves their divisor degrees,
so $w(aV)=w(V)$ for $a\in A$. Necessity follows from
Lemmas~\ref{lem:klyachko} and~\ref{lem:raycomponents}.
We prove sufficiency by linear algebra.

For any two subspaces $U,V\subseteq N_\CC$, counting the rays in their
union and intersection gives
\begin{equation}\label{eq:weightsupermodular}
 w(U)+w(V)\le w(U+V)+w(U\cap V).
\end{equation}
Let
\[
 \lambda=\max_{0\ne V\subseteq N_\CC}\frac{w(V)}{\dim V}.
\]
The maximum is attained among the finitely many subspaces spanned by
rays, and $\lambda\ge1$ because $w(N_\CC)=n$. If $U,V$ attain this
maximum, then \eqref{eq:weightsupermodular} and the bound
$w(W)\le\lambda\dim W$ imply that $U+V$ also attains it. They also
imply $w(U\cap V)=\lambda\dim(U\cap V)$, including when the
intersection is zero.
Every maximizing subspace is spanned by rays, since $\lambda>0$.
If $\lambda>1$, the sum of all maximizing subspaces is therefore a
nonzero proper maximizing subspace. It is spanned by rays and is
$A$-invariant, contrary to the hypothesis. Thus $\lambda=1$, and
$T_X$ is semistable by Lemma~\ref{lem:klyachko}.

Call a subspace $V$ \emph{tight} if $w(V)=\dim V$. Every nonzero tight
subspace is spanned by rays: replacing it by the span of its rays
would otherwise contradict $w(W)\le\dim W$.
The preceding argument shows that sums and intersections of tight
subspaces are tight. Moreover, if $U$ and $V$ are tight, equality
holds in \eqref{eq:weightsupermodular}. Since all ray weights are
positive, every ray in $U+V$ then lies in $U$ or in $V$.

There are finitely many nonzero tight subspaces. Let
$U_1,\ldots,U_q$ be the minimal ones under inclusion. Distinct $U_i$
have zero intersection, since an intersection is tight. Their sum is
in fact direct. To see this inductively, put $W=U_1+\cdots+U_{k-1}$.
The intersection $U_k\cap W$ is tight, so minimality implies that it
is either zero or all of $U_k$. In the latter case, choose a ray in
$U_k$. Repeated use of the preceding observation about rays in a sum shows that
this ray lies in some $U_i$ with $i<k$, contradicting
$U_i\cap U_k=0$. Thus
\[
 S=U_1\oplus\cdots\oplus U_q
\]
is tight, is spanned by rays, and is $A$-invariant, since $A$ permutes
the minimal nonzero tight subspaces.

If $S\ne N_\CC$, it contradicts the assumed strict inequality on
invariant subspaces. If $S=N_\CC$, every ray belongs to one of the
$U_i$, again by the same observation. Connectedness of the ray
matroid and Lemma~\ref{lem:raycomponents} force $q=1$. Hence
$U_1=N_\CC$ is minimal among nonzero tight subspaces, so there is no
proper nonzero tight subspace. All the inequalities in
Lemma~\ref{lem:klyachko} are therefore strict.
\end{proof}
